計量書誌学においては3年インパクトファクターや5年インパクトファクターなど\cite{Cla1}、特定の期間の（被）引用をもとにした指標が多く用いられる。
これは一面的には評価の公正性を考慮したものと考えられる。つまり、より古い論文ほど引用される機会は高くなり、この「論文の年齢」による効果を差し引くためである。
このような視点は雑誌の選定などにおいては重要なものである。一方で、長期に渡る引用の変化を把握することもまた重要である。第一に、前述の公正性における評価期間の設定の根拠となる。次に、短期的には捉えられない累積的な効果を捉えるには真に長期の観察が必要である。たとえば、引用という現象をより深く捉えるためには、単に（被）引用数に着目するだけでなく、引用関係とその変化などの構造的な面へのアプローチが望まれる。本研究では以上のような観点にたち、長期的な引用のダイナミクスを明らかにすることを試みる。
具体的には、1. 高被引用論文（定義は後述）の論文レベルでの長期的な被引用の構造、2. これらを掲載する雑誌の特徴（が見られるか）、3. これらの研究テーマ構造の長期的な変化（が見られるか）を、高被引用でない論文群と比較し分析する。
